% Pollis drill-down: the Pollis REPL, SQL and Julia at one prompt on the Shared Session, the two bridges between them,
% name := and @share, worked examples, and how the terminal reaches Julia (9 November 2026).
% For users who work with both languages, and for whoever maintains Pollis.
% Compile with: latexmk -pdf workflow_pollis_repl.tex
% Shared look for the one-page Pollis workflow handouts (same fonts and colours as the brochure).
\documentclass[10pt,a4paper]{article}

\usepackage[a4paper,left=17mm,right=17mm,top=14mm,bottom=16mm,footskip=8mm]{geometry}
\usepackage[T1]{fontenc}
\usepackage[default,semibold]{sourcesanspro}
\usepackage{microtype}
\usepackage{graphicx}
\usepackage[table]{xcolor}
\usepackage{tikz}
\usepackage[most]{tcolorbox}
\usepackage{fontawesome5}
\usepackage{enumitem}
\usepackage{fancyhdr}
\usepackage[hidelinks]{hyperref}

\graphicspath{{images/}}

% ---------------------------------------------------------------- colours (as in the brochure)
\definecolor{navy}{HTML}{1C3A5E}
\definecolor{pblue}{HTML}{3B78B5}
\definecolor{amber}{HTML}{F2B33D}
\definecolor{deepamber}{HTML}{D9902A}
\definecolor{pale}{HTML}{EEF3F9}
\definecolor{ink}{HTML}{2A2F36}
\definecolor{logonavy}{HTML}{173D63}
\definecolor{logodark}{HTML}{C98A1E}
\definecolor{logoamber}{HTML}{F2B446}
\definecolor{logocream}{HTML}{FEEFD2}
\definecolor{soft}{HTML}{5E6773}
\color{ink}

\setlength{\parindent}{0pt}
\setlength{\parskip}{3pt}
\setlist{nosep,leftmargin=*,itemsep=2pt}

\pagestyle{fancy}
\fancyhf{}
\renewcommand{\headrulewidth}{0pt}
\fancyfoot[L]{\footnotesize\color{soft}\textbf{Pollis}\enspace Bright Insights through Exploration, Modelling, Simulation and Optimisation}
\fancyfoot[R]{\footnotesize\color{soft}github.com/saragga/pollis/releases/latest}

% The Pollis logo, redrawn so it prints sharply at any size
\newcommand{\pollislogo}[1]{%
  \begin{tikzpicture}[x=#1/512,y=-#1/512,baseline=(current bounding box.center)]
    \fill[logonavy] (0,0) rectangle (512,512);
    \fill[logodark] (36,422) -- (104,256) -- (172,422) -- cycle;
    \fill[logoamber] (104,422) -- (203,180) -- (302,422) -- cycle;
    \fill[logocream] (203,422) -- (339,89) -- (476,422) -- cycle;
  \end{tikzpicture}}

% Top of the page: logo, name, kind of document (Workflow unless given, e.g. [Drill-Down]) and its subject;
% then the question and a one-line promise
\newcommand{\workflowhead}[4][Workflow]{%
  \thispagestyle{fancy}%
  \begin{minipage}[c]{13mm}\pollislogo{11mm}\end{minipage}%
  \begin{minipage}[c]{\dimexpr\linewidth-13mm\relax}
    {\color{navy}\Large\bfseries Pollis}\enspace{\color{soft}#1}\hfill
    {\color{deepamber}\small\bfseries\MakeUppercase{#2}}
  \end{minipage}\par\vspace{7mm}
  {\color{navy}\fontsize{22}{26}\selectfont\bfseries #3}\par\vspace{4pt}
  {\color{soft}\large #4}\par\vspace{5mm}}

% A menu path: \menu{Explore}{Access Data Libraries}... is written as \path{Explore \gt Access Data Libraries}
\newcommand{\gt}{\,{\color{deepamber}\bfseries ›}\,}
\newcommand{\menupath}[1]{{\footnotesize\bfseries\color{navy}#1}}

% A screenshot if images/<file> exists, otherwise an orange placeholder that says what to capture
\newlength{\shotheight}\setlength{\shotheight}{43mm}
\newcommand{\stepshot}[2]{%
  \IfFileExists{images/#1}{%
    \tcbox[enhanced, colback=black, colframe=black!70, boxrule=0.4pt, arc=2pt,
      left=0pt,right=0pt,top=0pt,bottom=0pt, boxsep=0pt, halign=center]{%
      \includegraphics[width=\dimexpr\linewidth-1pt\relax,height=\shotheight,keepaspectratio]{#1}}}{%
    \begin{tcolorbox}[enhanced, colback=white, colframe=orange!80, boxrule=0.6pt, arc=2pt,
      height=\shotheight, valign=center, halign=center, left=6pt, right=6pt]
      {\color{orange}\Large\faCamera}\par\smallskip
      {\color{orange}\footnotesize\bfseries Screenshot: #2}\par\smallskip
      {\color{orange!80!black}\scriptsize images/#1}
    \end{tcolorbox}}}

% One step: number, title, menu path, image file, what to capture, what it means
\newcommand{\step}[6]{%
  \begin{tcolorbox}
    {\color{deepamber}\Large\bfseries #1}\enspace{\color{navy}\large\bfseries #2}\par\vspace{1pt}
    \menupath{#3}\par\vspace{4pt}
    \stepshot{#4}{#5}\par\vspace{3pt}
    {\small #6}
  \end{tcolorbox}}

\newenvironment{steps}{%
  \begin{tcbraster}[raster columns=2, raster equal height=rows, raster column skip=4mm,
    raster row skip=4mm, enhanced, colback=pale, colframe=pale, boxrule=0pt, sharp corners,
    left=6pt,right=6pt,top=6pt,bottom=6pt, fontupper=\raggedright]}{\end{tcbraster}}

% Boxes for the bottom of the page
\newtcolorbox{pbox}[2][]{enhanced, colback=pale, colframe=amber, boxrule=0pt,
  leftrule=3pt, sharp corners, left=7pt, right=7pt, top=6pt, bottom=6pt,
  fontupper=\small\raggedright, before upper={{\color{navy}\normalsize\bfseries #2}\par\smallskip}, #1}
\newtcolorbox{nbox}[1][]{enhanced, colback=navy, colframe=navy, boxrule=0pt,
  sharp corners, left=8pt, right=8pt, top=6pt, bottom=6pt, coltext=white,
  fontupper=\small\raggedright, #1}

% A Next Steps item inside the bottom boxes
\newcommand{\nsitem}[2]{\item[\color{deepamber}\faChevronRight] \textbf{#1}: #2}

% The page of each heading, for Ask Pollis
% Pollis: the page of each heading, for Ask Pollis. A document that inputs this file, at the end of its preamble, writes
% <jobname>.labels beside its PDF, a line page|level|heading for each heading, the page counted from the first, the
% cover, as 1. compile_documents.jl copies it to ~/.pollis/documents, and ask_index.jl cuts the document into passages
% under its headings with it. \asklabel{1}{Title} marks a topic or a page, \asklabel{2}{Title} a heading inside it; the
% sections and subsections of a document that has them are marked as they go into its contents. A package vignette
% inputs it before its \documentclass (compile_vignette.jl), so the contents are taken over when the document begins,
% after the packages have taken them over.
\makeatletter
\newwrite\pollis@labels
\immediate\openout\pollis@labels=\jobname.labels
\newcommand{\asklabel}[2]{\write\pollis@labels{\the\ReadonlyShipoutCounter|#1|\unexpanded{#2}}}
\def\pollis@toc{toc}
\newcommand{\pollis@contentsline}[3]{%
  \def\pollis@file{#1}%
  \ifx\pollis@file\pollis@toc\asklabel{#2}{#3}\fi
  \pollis@addcontentsline{#1}{#2}{#3}}
% The headings of the New Methods style of the vignettes, \nmsection and \nmsub, which go into no contents
\def\pollis@asksection#1{\pollis@nmsection{#1}\asklabel{1}{#1}}
\def\pollis@asksub#1{\pollis@nmsub{#1}\asklabel{2}{#1}}
\AddToHook{begindocument/end}{%
  \let\pollis@addcontentsline\addcontentsline\let\addcontentsline\pollis@contentsline
  \ifdefined\nmsection\let\pollis@nmsection\nmsection\let\nmsection\pollis@asksection\fi
  \ifdefined\nmsub\let\pollis@nmsub\nmsub\let\nmsub\pollis@asksub\fi}
\makeatother



\usepackage{titlesec}
\usepackage{tabularx}
\usepackage{booktabs}
\usepackage{listings}

\fancyhf{}
\fancyfoot[L]{\footnotesize\color{soft}\textbf{Pollis}\enspace Drill-down: Pollis REPL}
\fancyfoot[R]{\footnotesize\color{soft}\thepage}

\titleformat{\section}{\color{navy}\large\bfseries}{\color{deepamber}\thesection}{0.6em}{}
\titleformat{\subsection}{\color{navy}\normalsize\bfseries}{\color{deepamber}\thesubsection}{0.5em}{}
\titlespacing*{\section}{0pt}{12pt}{4pt}
\titlespacing*{\subsection}{0pt}{8pt}{2pt}
\setlength{\parskip}{4pt}

\lstset{basicstyle=\ttfamily\footnotesize, columns=fullflexible, keepspaces=true, breaklines=true,
  backgroundcolor=\color{pale}, frame=leftline, rulecolor=\color{amber}, framerule=2pt, xleftmargin=6pt,
  framexleftmargin=6pt, aboveskip=4pt, belowskip=4pt, upquote=true, tabsize=2}

\usepackage{xurl}
\newcommand{\file}[1]{{\small\nolinkurl{#1}}}
\newcolumntype{L}{>{\raggedright\arraybackslash}X}

\begin{document}

\workflowhead[Drill-Down]{Pollis REPL}
  {SQL and Julia at one prompt}
  {The Pollis REPL runs SQL on the Shared Session, the database of the Julia REPL, and Julia code on the same Julia:
   how it tells them apart, how a table goes from one to the other and back, three worked examples, and what is under
   the surface.}

\begin{nbox}
{\color{amber}\bfseries Drill-down.}\enspace
A look under the surface of what the Guide describes in \emph{Data and Databases}, with its \emph{Four ways to use
SQL}. It is for anyone who moves between SQL and Julia in one piece of work, and for whoever maintains Pollis
(\S\ref{sec:machinery} is for them). Drill-down: Data Packages and Databases describes the databases themselves, and
how the \file{duckdb>} terminal and the notebooks reach the Shared Session. It describes the state of 9 November 2026.
\end{nbox}

\tableofcontents

\section{Three REPLs, one Julia}\label{sec:three}
Three items in the status bar, Julia REPL, Pollis REPL and DuckDB REPL, start the three REPLs of Pollis or show them;
the terminal's + menu offers the same three after its profiles. Each keeps to one job.

\begin{center}\small
\renewcommand{\arraystretch}{1.25}
\begin{tabularx}{\linewidth}{@{}p{0.15\linewidth}LL@{}}
\toprule
\textbf{Prompt} & \textbf{What you type} & \textbf{Where it runs, and on which database}\\
\midrule
\file{julia>} & Julia & the Julia REPL; \file{sql()} and \file{@share} reach the Shared Session\\
\file{pollis>} & SQL and Julia, told apart as you type & the same Julia: SQL on the Shared Session, Julia in the
  Julia REPL\\
\file{duckdb>} & SQL alone, and dot commands such as \file{.tables} & Pollis's own DuckDB program, without Julia, on a
  database of its own, which attaches the Shared Session as \file{SharedSession}\\
\bottomrule
\end{tabularx}
\end{center}
The Pollis REPL is a terminal of its own, beside the Julia REPL, so that \file{julia>} stays at hand. It runs on the
Julia of the Julia REPL, which it starts in the background if need be, and its item in the status bar is filled while
it is open. The Shared Session is the DuckDB of that Julia, in memory: its tables last until Julia exits, and are lost
when the Julia REPL restarts.

\section{One prompt, two languages}\label{sec:routing}
The Pollis REPL looks at the first word of what you type, after blanks and SQL comments.
\begin{itemize}
  \item \textbf{SQL}, when that word starts a statement of DuckDB and a blank follows it: \texttt{\small SELECT},
    \texttt{\small FROM}, \texttt{\small WITH}, \texttt{\small CREATE}, \texttt{\small INSERT},
    \texttt{\small DESCRIBE}, \texttt{\small SUMMARIZE}, \texttt{\small PIVOT} and the like, in any case; or when the
    input starts with \file{name :=}. SQL runs once it ends with \file{;}; until then, Enter starts a new line.
  \item \textbf{Julia}, anything else, such as \file{x = 1}, \file{using Statistics} or \file{@share t}. Julia runs
    once it is a whole expression, as at \file{julia>}: an open bracket or a \file{function} without its \file{end}
    asks for another line.
\end{itemize}
A name that is both, \file{describe}, \file{values} or \file{show}, is Julia when \file{(} follows it:
\texttt{\small DESCRIBE feeds;} is SQL, \file{describe(df)} Julia. \texttt{\small IMPORT}, \texttt{\small EXPORT} and
\texttt{\small BEGIN} are left to Julia.

SQL answers in the Pollis REPL, its result drawn to the terminal's width. Julia code runs in the Julia REPL, as Execute
Code does: the Julia REPL shows it with its output, its prints and its result in full, and the Pollis REPL writes the
result in one line, nothing for \file{nothing} or after a \file{;}, an error in red. Ctrl+C stops either: a query with
DuckDB's interrupt, Julia code as Ctrl+C in the Julia REPL does.

\section{The two bridges}\label{sec:bridges}
SQL sees the tables of the Shared Session; Julia sees its variables. Two bridges join them, one each way, and both
leave the table on both sides.

\begin{center}\small
\renewcommand{\arraystretch}{1.25}
\begin{tabularx}{\linewidth}{@{}p{0.36\linewidth}L@{}}
\toprule
\textbf{Typed} & \textbf{What it makes}\\
\midrule
\file{name := SELECT ...;} \hfill (SQL) & the table \file{name} of the Shared Session, and the Julia variable
  \file{name}, a \file{PollisDB.Table}\\
\file{@share name = table} \hfill (Julia) & the Julia variable \file{name}, as \file{=} would, and a copy of it as the
  table \file{name}\\
\file{@share name} \hfill (Julia) & the table \file{name} copied again, after a change in Julia\\
\file{sql("name", "SELECT ...")} \hfill (Julia) & what \file{name := SELECT ...} makes, from Julia code\\
\file{publish(table, "name")} \hfill (Julia) & what \file{@share} does, as a function\\
\bottomrule
\end{tabularx}
\end{center}
A Julia variable that has not been shared is not seen by SQL, and a table made by SQL without \file{:=} is not a Julia
variable. The bridges work the same way in the SQL and Julia cells of Julia notebooks, on the notebook's own database.

\subsection{From SQL to Julia: \texorpdfstring{\file{name :=}}{name :=} and \texorpdfstring{\file{PollisDB.Table}}{PollisDB.Table}}
\begin{lstlisting}
pollis> feeds := SELECT * FROM PollisRSS.feeds;
\end{lstlisting}
\file{PollisDB.pollis_query} runs it as \texttt{\small CREATE OR REPLACE TABLE "feeds" AS SELECT ...} on the Shared
Session, and binds the Julia variable \file{feeds} to a \file{PollisDB.Table}: the name of the table and its
connection, nothing more. Julia's multiple dispatch lets that one object serve both sides.
\begin{itemize}
  \item \textbf{To Julia}, it is a Tables.jl table: \file{feeds.title} is a column, and \file{length(feeds)} its
    number of rows, counted by DuckDB. Its rows are read at their first use and kept; a query of \file{sql} that
    changes tables, such as \texttt{\small INSERT} or \texttt{\small DROP}, makes every \file{Table} read its rows again.
  \item \textbf{To SQL}, it prints as its quoted name, so that an interpolation composes a query:
\begin{lstlisting}
sql("SELECT count(*) FROM $feeds")    # SELECT count(*) FROM "feeds"
\end{lstlisting}
  \item \textbf{In the Pollis REPL}, it is not drawn: \file{PollisDB.Kept} says its name, in Pollis's yellow, and its
    rows and columns, and under them \texttt{\small Browse feeds: Open in Table Viewer}, whose last words, blue and
    underlined, open the whole table in the Table Viewer with one click.
  \item \textbf{Under a notebook cell}, it reads only the rows it shows, the first 20 and the last 20, with the Julia
    type of each column.
\end{itemize}
A table that was dropped, or lost with its connection, says so when it is used, and to run again what made it.

\subsection{From Julia to SQL: \texorpdfstring{\file{@share}}{@share}}
\begin{lstlisting}
pollis> @share stock_returns = (ticker = ["AAPL", "MSFT"], returns = [0.12, 0.08])
\end{lstlisting}
\file{@share name = table} becomes \file{name = publish(table, "name")}, and \file{@share name} becomes
\file{publish(name, "name")}. Both are plain Julia, so they run anywhere Julia does: in the Pollis REPL, at
\file{julia>}, in code run from the editor, in notebooks, and inside functions and loops. \file{publish} registers the
table with DuckDB, copies it into the table \file{name} of the Shared Session, and returns it unchanged: Julia keeps
its own table, and a change to it reaches DuckDB when it is shared again. Any Tables.jl table is taken, a named tuple
of vectors as here, a \file{PollisDB.Table}, which is copied within DuckDB, or the table of another package; a value
that is not a table is refused, to be bound with \file{=}. A \file{missing} becomes \texttt{\small NULL}.

\section{Worked examples}\label{sec:examples}
Each example is typed at \file{pollis>}, a line or a statement at a time. The Julia lines run in the Julia REPL, which
shows their results in full; the Pollis REPL shows them in one line. They use only Julia's standard libraries and the
databases of Pollis.

\subsection{SQL, then Julia, then SQL}
Which topics of your feeds bring the most items? SQL takes the feeds, Julia scores them, and SQL sums the scores up.
\begin{lstlisting}
-- 1. SQL: the feeds, one row each, with the number of items each brought
pollis> feeds := SELECT feed, feed_title, topic, items FROM PollisRSS.feeds;
Table feeds, in the Shared Session and in Julia: 67 rows, 4 columns
Browse feeds: Open in Table Viewer

# 2. Julia: how far each feed is from the mean, in standard deviations
pollis> using Statistics
pollis> z = (feeds.items .- mean(feeds.items)) ./ std(feeds.items)

# 3. Back to SQL: the scores, with the columns that identify each feed
pollis> @share scores = (feed = feeds.feed, topic = feeds.topic, z = z)

-- 4. SQL again: group, filter and join as usual
pollis> SELECT topic, count(*) AS feeds, round(avg(z), 2) AS mean_z,
    ...>        count(*) FILTER (WHERE z > 1) AS busy_feeds
    ...> FROM scores GROUP BY topic ORDER BY mean_z DESC;

pollis> SELECT f.feed_title, round(s.z, 2) AS z
    ...> FROM scores s JOIN PollisRSS.feeds f USING (feed)
    ...> WHERE s.z > 1 ORDER BY s.z DESC;
\end{lstlisting}
Step 1 hands the result to Julia with \file{:=}; the numbers in its summary depend on your feeds. Step 2 is ordinary
Julia, and any package or function of yours can take its place: a z-score could be written in SQL too, and stands here
for what Julia does better, a model, a forecast or a simulation. Step 3 hands the result back with \file{@share}, and
step 4 reads \file{scores} by its name, beside a database of Pollis.

\subsection{Julia, then SQL, then Julia}
A Julia table of daily prices, simulated here, becomes returns in Julia; SQL subsets them and sums them up by month;
Julia fits a trend to what SQL sends back.
\begin{lstlisting}
# 1. Julia: a year of daily prices for three tickers, simulated with a fixed seed
pollis> using Random, Dates, Statistics
pollis> Random.seed!(1);
pollis> days = filter(d -> dayofweek(d) <= 5, Date(2025, 1, 1):Day(1):Date(2025, 12, 31));
pollis> prices = (ticker = repeat(["ALPHA", "BETA", "GAMMA"]; inner = length(days)),
    ...>           date = repeat(days; outer = 3),
    ...>           price = vcat([100 .* exp.(cumsum(s .* randn(length(days)))) for s in (0.010, 0.015, 0.020)]...));

# 2. Julia: the daily log returns, missing on each ticker's first day
pollis> starts = [i == 1 || prices.ticker[i] != prices.ticker[i - 1] for i in eachindex(prices.ticker)];
pollis> logret = [starts[i] ? missing : log(prices.price[i] / prices.price[i - 1]) for i in eachindex(prices.price)];

# 3. To SQL: the returns, a DuckDB table now, where missing is NULL
pollis> @share returns = (ticker = prices.ticker, date = prices.date, logret = logret);

-- 4. SQL: the volatility of each ticker over the year
pollis> SELECT ticker, count(*) AS days, round(stddev_samp(logret) * sqrt(252), 3) AS volatility
    ...> FROM returns WHERE logret IS NOT NULL GROUP BY ticker ORDER BY ticker;

-- 5. SQL: the second half of the year, by month, back to Julia
pollis> monthly := SELECT ticker, month(date) AS month, stddev_samp(logret) * sqrt(252) AS volatility
    ...> FROM returns WHERE logret IS NOT NULL AND date >= DATE '2025-07-01'
    ...> GROUP BY ticker, month ORDER BY ticker, month;
Table monthly, in the Shared Session and in Julia: 18 rows, 3 columns
Browse monthly: Open in Table Viewer

# 6. Julia: the trend of each ticker's monthly volatility, by least squares
pollis> X = [ones(6) 1:6]
pollis> trend = Dict(t => (X \ monthly.volatility[monthly.ticker .== t])[2] for t in unique(monthly.ticker))
\end{lstlisting}
The \file{;} at the end of a Julia line hides its result, as at \file{julia>}. Step 3 shares a table that Julia made
from nothing; step 4 shows a result without keeping it, and step 5 keeps one, which is a Julia table again in step 6:
\file{monthly.volatility} is a column, and \texttt{\small X \textbackslash{} y} Julia's least squares. Replace the simulated
prices with your own table.

\subsection{Your table beside the data of Pollis}
A list of journals in Julia, by their ISSNs, meets PollisJournals in one SQL join.
\begin{lstlisting}
# 1. Julia: the journals you aim at, and how many papers you have sent each
pollis> @share targets = (issn = ["0022-1082", "0304-405X", "0893-9454", "0022-1090"], submitted = [2, 1, 3, 0]);

-- 2. SQL: each journal's name, publisher and impact, from PollisJournals, back to Julia
pollis> ranked := SELECT j.title, j.publisher, j.h_index, j.mean_citedness_2yr AS citedness, t.submitted
    ...> FROM targets t
    ...> JOIN PollisJournals.journal_issns i ON i.issn = t.issn
    ...> JOIN PollisJournals.journals j ON j.id = i.id
    ...> ORDER BY j.h_index DESC;
Table ranked, in the Shared Session and in Julia: 4 rows, 5 columns
Browse ranked: Open in Table Viewer

# 3. Julia: the journal cited most in its last two years
pollis> ranked.title[argmax(ranked.citedness)]
\end{lstlisting}
\file{journal_issns} finds a journal by any of its ISSNs, print or electronic. The databases of Pollis are read-only;
your table sits beside them in the Shared Session, and the join reads both.

\section{Seeing and reaching the tables}\label{sec:seeing}
\begin{itemize}
  \item \textbf{Open in Table Viewer}, under a table that \file{:=} keeps, opens all of it in the Table Viewer, beside
    the Databases view, with one click, without the modifier other links of the terminal need.
  \item \textbf{Shared Session}, under User in the Databases view, lists the tables of the Shared Session, those of
    both bridges among them; click it to browse them. Its hover says how many clients of the window are attached to
    it and which of their queries are running.
  \item \textbf{\file{SharedSession.name}}: the \file{duckdb>} terminal and the notebooks of the window attach the
    Shared Session as \file{SharedSession}, so \texttt{\small SELECT * FROM SharedSession.scores} reads the table of
    the first example there, and \file{.tables} lists them. A query there may name a table of \file{SharedSession}
    once only; Drill-down: Data Packages and Databases gives the ways round.
  \item \textbf{History}, in the Databases view, lists the queries of the Pollis REPL, as \file{pollis>}, without
    their \file{name :=}, with their rows, how long they took and any error; Julia code is not listed.
\end{itemize}

\section{Good to know}\label{sec:know}
\begin{itemize}
  \item The Shared Session is in memory: its tables are lost when the Julia REPL restarts. Copy one to a file of your
    own to keep it, \texttt{\small COPY scores TO 'scores.parquet';}, or into a database of your own.
  \item One query runs at a time on the Shared Session; a second waits for the first. Julia code waits too, while the
    Julia REPL runs other code.
  \item \file{feeds := ...} again replaces the table and the Julia variable; \file{@share feeds} after a change in
    Julia replaces the table.
  \item The Pollis REPL and \file{julia>} share the variables of \file{Main}: a variable made in one is there in the
    other.
\end{itemize}

\section{Under the surface}\label{sec:machinery}
The Pollis REPL edits lines itself, with Node's readline (\file{extensions/language-julia/src/interactive/pollisTerminal.ts}),
and sends each whole input to Julia through a socket only you can reach, a named pipe on Windows, whose address and
token PollisDB writes to \file{~/.pollis/duckdb/served/pollis.<window>.json} when the REPL starts
(\file{PollisDB.serve_terminal}). The requests are lines, with the text they carry after them: \file{complete} asks
whether an input ends, \file{query} runs it, \file{browse} opens a table in the Table Viewer, \file{attached} asks what
the window's DuckDB clients have attached, and \file{interrupt} stops the query running. Julia answers each on threads
of its own, so that it answers while the REPL runs Julia code, and a lock keeps the Shared Session to one query at a
time.

\begin{itemize}
  \item \textbf{SQL or Julia}: \file{PollisDB.is_sql} decides, with \file{SQL_WORDS} and \file{ASSIGNMENT}; a Julia
    input is complete when \file{Base.parse_input_line} finds it so (\file{julia_complete}), an SQL one when it ends
    with \file{;} outside quotes and comments (\file{complete}).
  \item \textbf{SQL}: \file{PollisDB.terminal_query} runs it with \file{pollis_query}, logs it in
    \file{~/.pollis/duckdb/query_log.jsonl} with \file{"source": "pollis"}, which History reads, and sends back the
    result drawn to the terminal's width, or the error.
  \item \textbf{Julia}: Julia answers \file{julia}, and the terminal hands the input to the Julia REPL through the
    Julia extension's \file{executeInREPL}, as Execute Code does, without bringing the Julia REPL to the front. The
    REPL runs it in \file{Main}; Ctrl+C runs the Julia extension's interrupt.
  \item \textbf{The link}: Open in Table Viewer is a link of the terminal (\file{pollisTerminal.ts}, \file{browseLinks});
    the Julia extension asks Julia for the table through the socket, so nothing is typed into the REPL. It needs one
    click because the links of Pollis's own extensions do (\file{TerminalLink}, \file{oneClick}).
  \item \textbf{The keys before code}: the code Pollis sends to the Julia REPL, such as that of a page's Run button,
    comes after keys of no other use, ESC [ 6 9 \textasciitilde, which PollisDB binds in every mode of the REPL
    (\file{repl_keys}, \file{TO_JULIA}): to going to \file{julia>} from \file{pkg>} or \file{help>}, and in
    \file{julia>} to nothing; a Julia without PollisDB ignores them, as it does any keys ESC [ 6 and two more.
  \item \textbf{The status bar}: the three items keep to their names; the Julia REPL's threads and the database of the
    \file{duckdb>} terminal are in their tooltips.
\end{itemize}

\section{Where things are}

\begin{center}\small
\renewcommand{\arraystretch}{1.2}
\begin{tabularx}{\linewidth}{@{}p{0.47\linewidth}L@{}}
\toprule
\textbf{File} & \textbf{Role}\\
\midrule
\file{.../terminalserver/PollisDB.jl} & \file{sql}, \file{Table}, \file{Kept}, \file{publish}, \file{@share},
  \file{pollis_query}, \file{is_sql}, \file{serve_terminal}, \file{terminal_query}, \file{repl_keys}\\
\file{.../src/interactive/pollisTerminal.ts} & The Pollis REPL terminal, its status bar item, and the link Open in
  Table Viewer\\
\file{.../src/pollisReplConnection.ts} & Its connection to Julia\\
\file{.../src/interactive/repl.ts} & \file{executeInREPL}, which runs its Julia code in the Julia REPL\\
\file{.../terminal/browser/links/terminalLink.ts} & One-click links\\
\file{.../databases/browser/queryHistoryView.ts} & History\\
\bottomrule
\end{tabularx}
\end{center}

\end{document}
